\begin{frame}
	\frametitle{Modélisation géométrique}

	« dans le monde de tous les jours, la modélisation géométrique est utilisée
	dans de nombreux domaines » (illustrer: archi(revit), jv et animation
	(blender, godot, UE), pièces industrielles (catia, freecad))
\end{frame}

\begin{frame}
	\frametitle{Domaines d'applications}

	Il y a des domaines qui nécessitent de générer plusieurs variantes
	d'un même modèle (archéo, CAO…) (illustrer avec l'exemple de cardot, pièces mécanismes comme des boulons à tête hexa/octogonales)
\end{frame}

\begin{frame}
	\frametitle{Modélisation paramétrique}

	« dans ces travaux, nous nous focalisons sur la CAO et plus particulièrement
	la modélisation paramétrique. Parmi les avantages liés à ces outils se
	trouvent, d'une part, la possibilité de concevoir des pièces uniques, à
	l'aide d'opérations de modélisation géométrique, et d'autre part de
	regénérer ce modèle pour en créer des variations. »

	Focaliser le discours sur la CAO avec la
	modélisation paramétrique (avantages) (illustrer),

\end{frame}

\begin{frame}
	\frametitle{Processus de modélisation paramétrique}

	% Expliquer le fonctionnement d'un système de modélisation
	% paramétrique (application manuelle d'opérations, enregistrement du
	% processus de construction, rejeu)

	\begin{columns}[t]
		\begin{column}<1->{.33\textwidth}
			\begin{myblock}[Construction]
			  \vspace{.3\textheight}
			\end{myblock}
		\end{column}
		\begin{column}<2->{.33\textwidth}
			\begin{myblock}[Édition]
			  \vspace{.3\textheight}
			\end{myblock}
		\end{column}
		\begin{column}<3->{.33\textwidth}
			\begin{myblock}[Rejeu]
			  \vspace{.3\textheight}
			\end{myblock}
		\end{column}
	\end{columns}
\end{frame}


\begin{frame}[c,fragile]
	\frametitle{Problématique de la nomination persistante}

	% \begin{overlayarea}{\textwidth}{.7\textheight}
	\begin{columns}
		\begin{column}{.49\textwidth}
			\begin{minipage}[c][.5\textheight][b]{1.0\linewidth}
				\centering
				\includegraphics<1>[width=.7\textwidth]{1-square-thin}%
				\includegraphics<2>[width=.7\textwidth]{2-face-extrusion}%
				\includegraphics<3>[width=.7\textwidth]{4-face-triangulation}%
				\includegraphics<4>[width=.7\textwidth]{5-colour-face}%
				\includegraphics<5>[width=.7\textwidth]{1-square-thin}%
				\includegraphics<6>[width=.7\textwidth]{2-reeval-insert-vertex}%
				\includegraphics<7>[width=.7\textwidth]{3-reeval-face-extrusion}%
			\end{minipage}
		\end{column}
		\begin{column}{.49\textwidth}
			\centering
			\begin{minipage}[c][.7\textheight][t]{.7\linewidth}
				\vspace{.2\textheight}
				\begin{myblock}[\alt<5->{\vspace{1.3pt}Réévaluation}{Évaluation}]
					\ttfamily
					\only<1->{1-{\color{mydarkgreenink}square}({\color{orange}pos})}\\
					\only<1-4>{%
					\alt<2-4>{2-{\color{mydarkgreenink}extrusion}({\color{myblueink}f1},
					{\color{orange}vec})}
					{2-{\color{gray}extrusion}({\color{myblueink}f1},
					{\color{orange}vec})}
					\alt<3-4>{3-{\color{mydarkgreenink}triangulation}({\color{myblueink}f2})}
					{3-{\color{gray}triangulation}({\color{myblueink}f2})}\\
					\alt<4>{4-{\color{mydarkgreenink}colour}({\color{myblueink}f7})}
					{4-{\color{gray}colour}({\color{myblueink}f7})}%
					}%
					\only<5->{
					\alt<6->{{\color{red}ADD1-insert}({\color{myblueink}e1})}
					{ADD1-{\color{gray}insert({\color{myblueink}e1})}}
					\alt<7>{2-{\color{mydarkgreenink}extrusion}({\color{myblueink}f1},{\color{orange}vec})}
					{2-{\color{gray}extrusion}({\color{myblueink}f1},{\color{orange}vec})}
					\alt<7->{3-{\color{mydarkgreenink}triangulation}({\color{red}f2?})}
					{3-{\color{gray}triangulation}({\color{myblueink}f2})}\\
					\alt<7>{4-{\color{mydarkgreenink}colour}({\color{red}f7?})}
					{4-{\color{gray}colour}({\color{myblueink}f7})}%
					}
				\end{myblock}
			\end{minipage}
		\end{column}
	\end{columns}
	\begin{overlayarea}{\textwidth}{.1\textheight}
		\only<7>{%
			\begin{myblock}{}
				\color{red} \textrightarrow Nommage persistant nécessaire
			\end{myblock}
		}
	\end{overlayarea}
	\vfill
	\begin{overlayarea}{\textwidth}{.1\textheight}
		\begin{myblock}{}
			{\color{mydarkgreenink}Opérations} et paramètres {\color{myblueink}topologiques} et {\color{orange}géométriques}.
		\end{myblock}
	\end{overlayarea}
\end{frame}


\begin{frame}
	\frametitle{État de l'art pour la nomination persistante}

	\begin{itemize}
		\item utilisation d'un historique de noms
		\item nomination hétérogène des entités
	\end{itemize}
\end{frame}

\begin{frame}
	\frametitle{Limitations des approches existantes}

	\begin{itemize}
		\item absence de généricité dans les mécanismes de nomination
		      persistante (nommage différent en fonction de la dimension
		      de la cellule),
		\item modèles limités en dimension 3 (la plupart),
		\item édition limitée d'un processus constructif,
		\item (à compléter en fonction des contributions)
	\end{itemize}
\end{frame}

\begin{frame}
	\frametitle{Objectifs}
	Conception et développement d'un système de réévaluation permettant~:
	(autant d'objectifs que de limitations)
	\begin{itemize}
		\item une édition étendue des processus de construction (ajout,
		      suppression et réorganisation des opérations)~;
		      % \item de limiter  les coûts (calcul et mémoire) d'analyses et les risques
		      %       d'erreurs d'appariement~;
		\item de garantir la cohérence topologique des constructions.
	\end{itemize}
\end{frame}



% \begin{frame}[plain,noframenumbering]
% 	\frametitle{Plan}
% 	\tableofcontents
% \end{frame}




%%% Local Variables:
%%% mode: LaTeX
%%% TeX-master: "../../main"
%%% End:
